% ============================================================
% SOURCE MAP — DRAFT ONLY
% ============================================================
% Purpose: Working reference document mapping all dissertation
% sources by location (General / Movement I / II / III) and
% by tier (Keystone / Auxiliary / Contextual).
%
% Tier definitions:
%   [1] KEYSTONE  — Load-bearing for ST or for a movement's
%                   central argument. Cited at structural turns.
%                   The dissertation is fundamentally different
%                   without these.
%   [2] AUXILIARY — Develops, extends, or complicates a keystone
%                   claim in a specific section. Necessary but
%                   not structural. Thins the argument if absent.
%   [3] CONTEXTUAL — Establishes stakes, historical grounding,
%                    or a single specific claim. Often cited once.
%                    Provides the "so what" beyond the sonic/tech.
%
% UNPLACED: Sources not yet assigned to a section.
% Add \parencite{} keys as they are confirmed in references.bib.
%
% This file is DRAFT ONLY. It is guarded by \ifdraft in
% main.tex and will NOT appear in the final dissertation.
% ============================================================

% ------------------------------------------------------------
% HOW TO USE THIS MAP
% ------------------------------------------------------------

\section*{How to Use This Map}

This is a working reference for tracking where sources live in the dissertation, how much weight they carry, and how they relate to each other. It is organized in two dimensions:

\textbf{Location}: General (cross-cutting, not owned by a single movement), or specific to Movement I (Ground), II (Hold), or III (Touch). Some sources appear in multiple locations — they are listed primarily where they do the most structural work.

\textbf{Tier}: 1 (Keystone), 2 (Auxiliary), 3 (Contextual). Tier assignments are provisional and should be updated as the writing develops. A source that begins as contextual may become keystone if the argument evolves to depend on it.

Use this document when writing to quickly locate which sources are available in a given section, and which gaps still need filling. The \textit{Unplaced} section at the end flags sources that need to be assigned — or confirmed as out of scope.

\bigskip

% ------------------------------------------------------------
% GENERAL (CROSS-CUTTING)
% Sources that operate across the full score — not owned
% by any single movement. These are the theoretical spine.
% ------------------------------------------------------------

\section*{General: Cross-Movement Sources}
\addcontentsline{toc}{section}{General: Cross-Movement Sources}

These sources anchor ST itself, the research-creation methodology, and the sonic-epistemological claims that all three movements share. They appear first in the Composer's Note and recur throughout.

\medskip

\subsection*{[1] Keystone — General}

\begin{itemize}

    \item \textbf{Agre, Philip — ``Toward a Critical Technical Practice''} \parencite{Agre_1998}\\
    The methodological spine. Defines the practitioner who turns a discipline's tools on the discipline itself, producing knowledge through modification rather than use. Every instrument in this dissertation is an instance of CTP. Referenced explicitly in Movement III (HairSynth) and operative implicitly everywhere.\\
    \textit{Relation}: Pairs with Sengers et al.\ (Reflective Design) as the more technically-grounded / more HCI-grounded versions of the same claim. Agre is the harder, more self-implicating version.\\
    \admargin{Key mismatch: advisory-outline.tex uses \texttt{Agre\_1997}; bib key is \texttt{Agre\_1998}. Update advisory-outline.tex citations to \texttt{Agre\_1998}.}

    \item \textbf{Wynter, Sylvia — ``Unsettling the Coloniality of Being''} \parencite{Wynter_2003}\\
    The epistemic ground for ST's claim that modifying tools is modifying the embedded assumptions about who counts as a knowing subject. The over-representation of ``Man'' is the target; ST is one form of unsettling it.\\
    \textit{Relation}: Pairs with Wynter 2015 (``The Ceremony Found'') — 2003 names the problem, 2015 names the mode of response.

    \item \textbf{Wynter, Sylvia — ``The Ceremony Found''} \parencite{Wynter_2015}\\
    The autopoetic act of re-envisioning the human beyond Western epistemic closure. ST as a sonic ceremony-finding. Primary citation in the Composer's Note positionality section.

    \item \textbf{Benjamin, Ruha — \textit{Race After Technology}} \parencite{Benjamin_2019a}\\
    ``The road to inequity is paved with technical fixes.'' Surface-level interventions leave the epistemic foundation untouched. Appears in multiple movements as the counter-argument ST is built against.

    \item \textbf{Barad, Karen — \textit{Meeting the Universe Halfway}} \parencite{Barad_2007}\\
    Intra-action and apparatus. Each instrument is not a neutral tool applied to a cultural problem — the instrument, the cultural encoding, and the sonic output emerge together. Appears in all three ``Instrument Built'' sections. Foundational for how ST understands the relationship between making and knowledge.

    \item \textbf{Grimshaw, Mark \& Garner, Tom — \textit{Sonic Virtuality}}\\
    Transduction and refraction. Sound as a field of potential, shaped by the conditions of its mediation. Provides the sonic-philosophical vocabulary for what all three instruments do. The drone as pre-unsettled condition. Appears in Composer's Note and returns in Movement III.\\
    \textit{Key terms}: transduction, refraction, sonic virtuality, sonic mediation.\\
    \textit{Note}: Confirm \texttt{\\parencite} key in references.bib.

    \item \textbf{Ihde, Don — \textit{Technology and the Lifeworld}} \parencite{Ihde_1990}\\
    The hermeneutic relation: instruments shape what can be perceived by the design decisions embedded in them. The thread that connects all three sonic sculptures. Appears in Composer's Note, Movement I, and Movement II. Every sensor in this dissertation is an Ihde hermeneutic relation.

\end{itemize}

\subsection*{[2] Auxiliary — General}

\begin{itemize}

    \item \textbf{Manning, Erin \& Massumi, Brian — \textit{Thought in the Act}} \parencite{Manning_Massumi_2014}\\
    Thinking-feeling; the inseparability of conceptual and bodily engagement in making. Supports the research-creation claim that theory and practice are not sequential but simultaneous. Appears in ``Making the Score'' (Composer's Note).\\
    \textit{Relation}: Extends Loveless methodologically; grounds the embodied dimension of ST.

    \item \textbf{Loveless, Natalie — \textit{How to Make Art at the End of the World}} \parencite{Loveless_2019}\\
    Research-creation as inquiry driven by practice. Methodological grounding for the dissertation's form. Appears in ``Making the Score'' (Composer's Note).\\
    \textit{Relation}: Pairs with Manning \& Massumi; together they make the methodological case.

    \item \textbf{Franinović, Karmen \& Serafin, Stefania — \textit{Sonic Interaction Design}} \parencite{Franinovic_Serafin_2013}\\
    Interaction design frameworks applied to body-driven sonic systems. Appears in all three ``Instrument Built'' sections — providing the technical design vocabulary Barad's intra-action requires in practice. The PebbleBox is cited as direct precedent for Movement I.\\
    \textit{Relation}: Pairs with Barad in every ``Instrument Built'' section — Barad names what the apparatus does; Franinović \& Serafin names how you design it.

    \item \textbf{Eshun, Kodwo — \textit{More Brilliant Than the Sun}} \parencite{Eshun_1999}\\
    Sonic fiction as speculative apparatus; temporal reconfiguration through sound. Appears in Composer's Note (``Trombone in Third Position'') and is the primary theoretical lens for Movement II (``Entry: The Question''). Also returns in Log 3 (Temporal Escape).\\
    \textit{Relation}: Pairs with Womack (Afrofuturism) in Movement II — Eshun is the sonic/speculative frame; Womack provides the broader cultural-historical frame.

    \item \textbf{Eastman, Julius — \textit{Femenine}}\\
    The structural model: drone as perceptual constant, developmental cell as generative unit, duration without resolution. Cited in ``Making the Score.'' Not a theoretical source — a compositional precedent and an argument for the form.\\
    \textit{Note}: Confirm citation format and key. This may need to be a score citation, not a text.

\end{itemize}

\subsection*{[3] Contextual — General}

\begin{itemize}

    \item \textbf{Marks, Laura — \textit{The Skin of the Film}} \parencite{Marks_2000}\\
    Haptic visuality; the skin as an epistemological organ. Primarily serves Movement III (Section 3.0, Haptic Epistemology), but the broader claim — that touch is a way of knowing that precedes and exceeds vision — runs as background assumption across all three works. Listed here as General because the epistemological claim is broader than Movement III alone.

\end{itemize}

\bigskip

% ------------------------------------------------------------
% MOVEMENT I: GROUND
% Buffalo Soldier — earth, rock, substrate, piezo, tactility
% ST Domain: Material Memory (MM)
% ------------------------------------------------------------

\section*{Movement I: Ground (\textit{Buffalo Soldier})}
\addcontentsline{toc}{section}{Movement I: Ground}

\subsection*{[1] Keystone — Movement I}

\begin{itemize}

    \item \textbf{Sharpe, Christina — \textit{In the Wake}} \parencite{Sharpe_2016}\\
    \textit{Note}: Sharpe appears in Movement II (Log 1: The Wake) as the primary theoretical lens. Listed there. However, the conceptual infrastructure of the ship, the hold, and the weather as methodology is also operative in Movement I's thinking about what the earth holds. Consider whether a secondary citation in Ground's synthesis (``Deep Layer: What Ground Hears'') would be warranted.

    \item \textbf{McKittrick, Katherine — \textit{Demonic Grounds}} \parencite{McKittrick_2006}\\
    Black geographies; land as site of contested demarcation, extraction, and burial. The foundational text for reading the substrate box as archive. Appears in ``First Layer: The Ground as Archive.''

    \item \textbf{Hartman, Saidiya — ``Venus in Two Acts'' / Critical Fabulation} \parencite{Hartman_2008}\\
    Filling archival silence through speculative method. The encounter between soldier and buffalo has no archive; granular synthesis enacts the fabulation. Central to ``Second Layer: The Encounter That Has No Archive.'' Without Hartman, the speculative compositional method has no methodological name.

\end{itemize}

\subsection*{[2] Auxiliary — Movement I}

\begin{itemize}

    \item \textbf{Bennett, Jane — \textit{Vibrant Matter}} \parencite{Bennett_2010}\\
    Thing-power; nonhuman agency. Rocks, soil, and sand co-author the sonic output — the material makes decisions the patch did not anticipate. Grounds the claim that the substrate box is a collaborator, not a container. Appears in ``First Layer.''\\
    \textit{Relation}: Pairs with Barad (General) — Bennett names the thing-power of individual materials; Barad names what the full apparatus (materials + sensors + code + cultural logic) produces together.

    \item \textbf{Glissant, Édouard — \textit{Poetics of Relation}} \parencite{Glissant_1997}\\
    Opacity as a right; the refusal of full legibility to dominant systems of knowledge. Grounds the intentional sonic ambiguity of the piezo signal — flesh, stone, soil, breath are not meant to be cleanly distinguished. Appears in ``Third Layer: The Right to Not Be Legible.''

\end{itemize}

\subsection*{[3] Contextual — Movement I}

\begin{itemize}

    \item \textbf{Estes, Nick — \textit{Our History Is the Future}} \parencite{Estes_2019}\\
    The deliberate near-extinction of the bison as colonial weapon; the triangulated histories in the name ``Buffalo Soldier.'' Provides the historical grounding for the Unsettling in ``Surface: The Desire and the Unsettling.'' Essential stakes-setting, but the argument does not turn on Estes — it turns on what the instrument does with the history Estes documents.

\end{itemize}

\bigskip

% ------------------------------------------------------------
% MOVEMENT II: HOLD (PRIMARY)
% Did Sun Ra Fear Cryosleep? — water, hydrophone, dissolution
% ST Domain: Fugitive Speculation (FS)
% ------------------------------------------------------------

\section*{Movement II: Hold (\textit{Did Sun Ra Fear Cryosleep?}) --- PRIMARY}
\addcontentsline{toc}{section}{Movement II: Hold (PRIMARY)}

\adnote[title={Primary Movement}]{Movement II receives the deepest theoretical development. Sources here have been or are being actively engaged in the writing from April 2026 onward. All source placements in this section should be considered load-bearing.}

\subsection*{[1] Keystone — Movement II}

\begin{itemize}

    \item \textbf{Sharpe, Christina — \textit{In the Wake}} \parencite{Sharpe_2016}\\
    The ship, the hold, the wake, the weather; wake work as methodology; the Atlantic as both grave and archive. The primary theoretical lens for ``Log 1: The Wake.'' Also gives the movement its central claim: the default signal processing that ``cleans'' hydrophone recordings erases the residue Sharpe says must be tended to. The modification at the processing chain is an act of wake work.

    \item \textbf{Eshun, Kodwo — \textit{More Brilliant Than the Sun}} \parencite{Eshun_1999}\\
    Already listed General (Auxiliary). Keystone here specifically: the speculative question (``Did Sun Ra fear cryosleep?'') is the research instrument Eshun names — it opens rather than closes. Re-listed at keystone level for Movement II because Entry: The Question is the architectural center of this movement.

    \item \textbf{Philip, M.\ NourbeSe — \textit{Zong!}} \parencite{Philip_2008}\\
    Language fragmented and drowned. Philip does with text what the hydrophone does with water: makes the Atlantic archive audible by fragmenting the dominant narrative. The formal parallel is the argument. Appears in ``Log 2: What the Water Holds.'' The processing chain design follows directly from Philip's formal logic.

    \item \textbf{Phillips, C.\ S.\ / Black Quantum Futurism} \parencite{Phillips_2015}\\
    Black non-linear temporality; cryosleep as temporal technology. The theoretical engine for ``Log 3: Temporal Escape.'' Without Black Quantum Futurism, the move from Middle Passage to interstellar transit is a metaphor; with it, it is an argument about time, survival, and the colonial timeline.\\
    \textit{Note}: Confirm full citation — is this Rasheedah Phillips, \textit{Black Quantum Futurism: Theory \& Practice}? Verify key \texttt{Phillips\_2015} in references.bib.

\end{itemize}

\subsection*{[2] Auxiliary — Movement II}

\begin{itemize}

    \item \textbf{Neimanis, Astrida — \textit{Bodies of Water}} \parencite{Neimanis_2017}\\
    Trans-corporeality; we are bodies of water, not bodies in it. Dissolution as homecoming, not only loss. Reframes what the hydrophone is listening for in ``Log 2,'' and gives the processing chain a positive direction (toward transformation) rather than only a negative one (against erasure).\\
    \textit{Relation}: Pairs with Sharpe — Sharpe holds the grief; Neimanis holds the transformation. Both are needed.

    \item \textbf{Butler, Octavia — \textit{Xenogenesis / Lilith's Brood}} \parencite{Butler_1987}\\
    Lilith wakes from suspension into an alien ship; the Black body simultaneously preserved and violated by what holds it. The literary/speculative precedent for the cryo-pod motif in ``Log 3.'' Grounds the claim that Black speculative fiction has already thought through what the visual archive erases.\\
    \textit{Relation}: Pairs with Womack (Afrofuturism) as the creative practice to Womack's cultural-historical frame.

    \item \textbf{Womack, Ytasha — \textit{Afrofuturism}} \parencite{Womack_2013}\\
    Technology, liberation, and the construction of Black futurity; Sun Ra's mythology as Afrofuturist practice. Provides the cultural-historical grounding for the speculative question. Appears in ``Entry: The Question.''\\
    \textit{Relation}: Pairs with Eshun — Eshun is the sonic-speculative apparatus; Womack is the cultural-historical frame.

\end{itemize}

\bigskip

% ------------------------------------------------------------
% MOVEMENT III: TOUCH
% Don't Play with My Hair — hair, gesture, CV, HairSynth
% ST Domain: Sensory Worlding (SW)
% ------------------------------------------------------------

\section*{Movement III: Touch (\textit{Don't Play with My Hair})}
\addcontentsline{toc}{section}{Movement III: Touch}

\subsection*{[1] Keystone — Movement III}

\begin{itemize}

    \item \textbf{Spillers, Hortense — ``Mama's Baby, Papa's Maybe''} \parencite{Spillers_1987}\\
    The distinction between body and flesh; the flesh as the site of racial inscription; the uninvited touch as a claim on the flesh. The HairSynth's central argument depends on this distinction: the reclamation occurs at the level of flesh, not body — knowledge and sound generated on the flesh's own terms. Appears in ``Before the Manual Begins'' and returns in Section 1.0.

    \item \textbf{Buolamwini, Joy — \textit{Gender Shades}} \parencite{Buolamwini_2018}\\
    CV systems fail to detect Black hair because training data encoded whiteness as norm; Type 4 hair is structurally absent from most datasets. The technical unsettling that the HairSynth is built to confront. Appears in ``Section 2.0: The Pipeline.'' Without Buolamwini, the CV failure is a technical limitation; with her, it is an argument the algorithm is making.

    \item \textbf{Mercer, Kobena — ``Black Hair/Style Politics''} \parencite{Mercer_1987}\\
    Hair as signifying practice; aesthetics and politics as inseparable. The HairSynth makes this signifying practice computationally generative — the sign generates sound. Appears in ``Section 1.0: The Body as Signifying Practice.''

\end{itemize}

\subsection*{[2] Auxiliary — Movement III}

\begin{itemize}

    \item \textbf{Marks, Laura — \textit{The Skin of the Film}} \parencite{Marks_2000}\\
    Haptic visuality; the skin as epistemological organ; touch as a way of knowing that precedes and exceeds vision. Names why the embodied data-mapping process (adjusting while listening until the output feels right) is methodology rather than subjectivity. Appears in ``Section 3.0: Haptic Epistemology.''\\
    \textit{Relation}: Also listed General (Contextual) for the broader epistemological claim.

    \item \textbf{Kimmerer, Robin Wall — \textit{Braiding Sweetgrass}} \parencite{Kimmerer_2013}\\
    Braiding as a practice of attending to material with patience, relationship, and intergenerational knowledge. Provides a framework for reading the intimate gestures the HairSynth is built to receive. Appears in ``Section 4.0: Braiding.''\\
    \textit{Note}: \admargin{Distinct traditions — handle carefully. Do not collapse Indigenous and Black hair care epistemologies. The connection is formal (braiding as relational practice) not equivalent.}

\end{itemize}

\subsection*{[3] Contextual — Movement III}

\begin{itemize}

    \item \textbf{Benjamin, Ruha — \textit{Race After Technology}} \parencite{Benjamin_2019a}\\
    Also listed General (Keystone). Appears here specifically in Movement III (Section 2.0) as the interpretive frame for the CV failure: the algorithm's failure is the argument of the algorithm. Single citation, but structurally important for naming what the pipeline is doing rather than just failing to do.

\end{itemize}

\bigskip

% ------------------------------------------------------------
% UNPLACED / TO BE LOCATED
% Sources not yet assigned — need home or scope decision
% ------------------------------------------------------------

\section*{Unplaced Sources — To Be Located}
\addcontentsline{toc}{section}{Unplaced Sources}

\adnote[title={Action Required}]{The sources below need to be placed or explicitly scoped out. For each, the likely location and tier are suggested, but the final assignment depends on the writing. ``Sonic Writing'' and ``Sonic Thinking'' are particularly time-sensitive given their relevance to the Composer's Note.}

\begin{itemize}

    \item \textbf{Magnusson, Thor — \textit{Sonic Writing: Technologies of Material, Symbolic, and Signal Inscriptions}} (2019) \parencite{Magnusson_2019}\\
    \textit{Suggested location}: General (Composer's Note, ``Making the Score'' / ``A Trombone in Third Position''), with possible return in all three ``Instrument Built'' sections.\\
    \textit{Suggested tier}: [2] Auxiliary — General.\\
    \textit{Rationale}: Magnusson theorizes notation as a cognitive tool, instrument design as a form of writing, and the score as an epistemological artifact — which is exactly what this dissertation's form argues. The trombone story is already a story about what the instrument writes into the player. Magnusson gives that a technical vocabulary. He extends Agre into the specifically sonic domain: if CTP is the practitioner turning tools on themselves, Sonic Writing is how instruments themselves encode epistemological commitments. A bridge between the music-technical side and the critical-theoretical side of the dissertation. Citation key \texttt{Magnusson\_2019} added to references.bib.

    \item \textbf{Herzogenrath, Bernd (ed.) — \textit{Sonic Thinking: A Media Philosophical Approach}} (2017) \parencite{Herzogenrath_2017}\\
    \textit{Suggested location}: General (Composer's Note), ``The Drone'' or ``How to Hear This.''\\
    \textit{Suggested tier}: [2] Auxiliary — General.\\
    \textit{Rationale}: Positions sound as a mode of thought across disciplines — not just a medium but an epistemological practice. Situates the Composer's Note's opening claim (sound arrives before the instrument shapes it; the drone is the pre-unsettled condition) in a broader philosophical context. Functions as the opening theoretical move that Grimshaw \& Garner then specify technically. Together they'd make: sound as thought (Herzogenrath) $\to$ sonic transduction and refraction as mechanism (Grimshaw \& Garner) $\to$ instrument design as epistemological act (Magnusson) $\to$ CTP as methodology (Agre). Citation key \texttt{Herzogenrath\_2017} added to references.bib.

    \item \textbf{Sengers, Phoebe et al.\ — ``Reflective Design''} (2005)\\
    \textit{Suggested location}: General (Composer's Note, ``Making the Score'') or Movement III (Section 5.0: Instrument Built).\\
    \textit{Suggested tier}: [2] Auxiliary — General.\\
    \textit{Rationale}: The HCI-side complement to Agre's CTP. Reflective design asks designers to surface and interrogate the values embedded in their systems — directly applicable to HairSynth and the substrate box. Where Agre comes from AI/computer science and is harder and more self-implicating, Sengers comes from HCI and is more process-oriented. Together they bracket the scope of what CTP means in practice.\\
    \textit{Note}: Not currently in references.bib — add citation key \texttt{Sengers\_2005} or confirm.

\end{itemize}

\bigskip

% ------------------------------------------------------------
% QUICK REFERENCE: ALL SOURCES BY TIER
% ------------------------------------------------------------

\section*{Quick Reference: All Sources by Tier}
\addcontentsline{toc}{section}{Quick Reference by Tier}

\subsection*{[1] Keystones}

\begin{itemize}[noitemsep]
    \item Agre (CTP) — General
    \item Wynter 2003 (Unsettling Coloniality) — General
    \item Wynter 2015 (Ceremony Found) — General
    \item Benjamin (Race After Technology) — General
    \item Barad (Meeting the Universe Halfway) — General
    \item Grimshaw \& Garner (Sonic Virtuality) — General
    \item Ihde (Technology and the Lifeworld) — General
    \item McKittrick (Demonic Grounds) — Mvmt I
    \item Hartman (Critical Fabulation) — Mvmt I
    \item Sharpe (In the Wake) — Mvmt II (primary)
    \item Eshun (More Brilliant) — General + Mvmt II
    \item Philip (\textit{Zong!}) — Mvmt II
    \item Phillips / Black Quantum Futurism — Mvmt II
    \item Spillers (Mama's Baby) — Mvmt III
    \item Buolamwini (Gender Shades) — Mvmt III
    \item Mercer (Black Hair) — Mvmt III
\end{itemize}

\subsection*{[2] Auxiliaries}

\begin{itemize}[noitemsep]
    \item Manning \& Massumi (Thought in the Act) — General
    \item Loveless (Research-Creation) — General
    \item Franinović \& Serafin (Sonic Interaction Design) — General (all ``Instrument Built'')
    \item Eastman (\textit{Femenine}) — General (structural precedent)
    \item Magnusson (\textit{Sonic Writing}) — General [UNPLACED, in bib]
    \item Herzogenrath (ed.) (\textit{Sonic Thinking}) — General [UNPLACED, in bib]
    \item Sengers et al.\ (Reflective Design) — General [UNPLACED, in bib]
    \item Bennett (Vibrant Matter) — Mvmt I
    \item Glissant (Poetics of Relation) — Mvmt I
    \item Neimanis (Bodies of Water) — Mvmt II
    \item Butler (\textit{Lilith's Brood}) — Mvmt II
    \item Womack (Afrofuturism) — Mvmt II
    \item Marks (Skin of the Film) — General + Mvmt III
    \item Kimmerer (Braiding Sweetgrass) — Mvmt III
\end{itemize}

\subsection*{[3] Contextuals}

\begin{itemize}[noitemsep]
    \item Estes (\textit{Our History Is the Future}) — Mvmt I
    \item Benjamin (Race After Technology) — Mvmt III (Section 2.0 specific)
\end{itemize}
